\chapter{2009年重庆卷（理）}

\section{单选题}

\begin{problem}
直线 \( y = x + 1 \) 与圆 \( x^{2} + y^{2} = 1 \) 的位置关系是（\quad）

\choices
    {相切}
    {相交但直线不过圆心}
    {直线过圆心}
    {相离}
\end{problem}

\begin{problem}
已知复数 \(x\) 的实部为 \(-1\)，虚部为 \(2\)，则 \(\frac{5\i}{4} =\)（\quad）

\choices
    {2 - i}
    {\(2 + \i\)}
    {\(-2\) - i}
    {\(-2 + \i\)}
\end{problem}

\begin{problem}
\(\left(x^{2} + \frac{2}{x}\right)^{8}\) 的展开式中 \(x^{4}\) 的系数是（\quad）

\choices
    {16}
    {70}
    {560}
    {1120}
\end{problem}

\begin{problem}
已知 \(|\pmb{a}| = 1, |\pmb{b}| = 6, \pmb{a} \cdot (\pmb{b} - \pmb{a}) = 2\)，则向量 \(\pmb{a}\) 与 \(\pmb{b}\) 的夹角是（\quad）

\choices
    {\(\frac{\pi}{6}\)}
    {\(\frac{\pi}{4}\)}
    {\(\frac{\pi}{3}\)}
    {\(\frac{\pi}{2}\)}
\end{problem}

\begin{problem}
不等式 \( |x + 3| - |x - 1| \leqslant a^2 - 3a \) 对任意实数 \( x \) 组成元，则实数 \( a \) 的取值范围为（\quad）

\choices
    {\((-\infty, -1] \cup [4, +\infty)\)}
    {\((-\infty, -2] \cup [5, +\infty)\)}
    {\( [1,2] \)}
    {\((-\infty, 1] \cup [2, +\infty)\)}
\end{problem}

\begin{problem}
锅中煮有芝麻馅汤圈 6 个，花生馅汤圈 5 个，豆沙馅汤圈 4 个，这三种汤圈的外部特征完全相同. 从中任意四取 4 个汤圈，则每种汤圈都至少取到 1 个的概率为（\quad）

\choices
    {\(\frac{3}{91}\)}
    {\(\frac{25}{91}\)}
    {\(\frac{48}{91}\)}
    {\(\frac{60}{91}\)}
\end{problem}

\begin{problem}
设 \(\triangle ABC\) 的三个内角 \(A, B, C\)，向量 \(m = (\sqrt{3} \sin A, \sin B)\)，\(n = (\cos B, \sqrt{3} \cos A)\)，若 \(m \cdot n = 1 + \cos (A + B)\)，则 \(C =\)（\quad）

\choices
    {\(\frac{5}{6}\)}
    {\(\frac{\pi}{3}\)}
    {\(\frac{2\pi}{3}\)}
    {\(\frac{5\pi}{6}\)}
\end{problem}

\begin{problem}
已知 \(\lim_{x\to \infty}\left(\frac{2x^2}{x + 1} -ax - b\right) = 2\) ，其中 \(a,b\in \R\) ，则 \(a - b\) 的值为（\quad）

\choices
    {\(-6\)}
    {\(-2\)}
    {2}
    {6}
\end{problem}

\begin{problem}
已知二面角 \(\alpha - l - \beta\) 的大小为 \(50^{\circ}\)，\(P\) 为空间中任意一点，则过点 \(P\) 且与平面 \(\alpha\) 和平面 \(\beta\) 所成的角都是 \(25^{\circ}\) 的直线的条数为（\quad）

\choices
    {2}
    {3}
    {4}
    {5}
\end{problem}

\begin{problem}
已知以 \(T = 4\) 为周期的函数 \(f(x) = \begin{cases} m\sqrt{1 - x^2}, & x \in (-1, 1], \\ 1 - |x - 2|, & x \in (1, 3]. \end{cases}\) 其中 \(m > 0\). 若方程 \(3f(x) = x\) 恰有5个实数解，则 \(m\) 的取值范围为（\quad）

\choices
    {\(\left(\frac{\sqrt{15}}{3}, \frac{8}{3}\right)\)}
    {\(\left(\frac{\sqrt{15}}{3}, \sqrt{7}\right)\)}
    {\(\left(\frac{4}{3}, \frac{8}{3}\right)\)}
    {\(\left(\frac{4}{3}, \sqrt{7}\right)\)}
\end{problem}

\begin{problem}
若 \(A = \{x \in \R \mid |x| < 3\}\)，\(B = \{x \in \R \mid 2^x > 1\}\)，则 \(A \cap B = \fillinblank{}\).
\end{problem}

\begin{problem}
若 \( f(x) = \frac{1}{2^x - 1} + a \) 是奇函数，则 \( a = \) .
\end{problem}

\begin{problem}
将 4 名大学生分配到 3 个乡镇去当村官，每个乡镇至少一名，则不同的分配方案有 种. (用数字作答)
\end{problem}

\begin{problem}
设 \(a_1 = 2, a_{n+1} = \frac{2}{a_{n+1}}, b_n = \left|\frac{a_n + 2}{a_n - 1}\right|, n \in \mathbf{N}^*\)，则数列 \(\{b_n\}\) 的通项 \(b_n = \) .
\end{problem}

\begin{problem}
已知双曲线 \(\frac{x^2}{a^2} - \frac{y^2}{b^2} = 1 (a > 0, b > 0)\) 的左、右焦点分别为 \(F_1(-c, 0)\)，\(F_2(c, 0)\). 若双曲线上存在一点 \(P\) 使 \(\frac{\sin \angle P_1 F_2}{b \sin \angle P_2 F_1} = \frac{a}{c}\)，则该双曲线的离心率的取值范围为 .
\end{problem}

\section{解答题}

\begin{problem}
设函数 \( f(x) = \sin \left(\frac{\pi}{3} x - \frac{\pi}{6}\right) - 2\cos^2 \frac{\pi}{8} x + 1 \). (1) 求 \( f(x) \) 的最小正周期； (2) 若函数 \( y = g(x) \) 与 \( y = f(x) \) 的图象关于直线 \( x = 1 \) 对称，求当 \( x \in \left[0, \frac{4}{3}\right] \) 时 \( y = g(x) \) 的最大值.
\end{problem}

\begin{problem}
设函数 \( f(x) = ax^2 + bx + k (k > 0) \) 在 \( x = 0 \) 处取得极值，且曲线 \( y = f(x) \) 在点 \( (1, f(1)) \) 处的切线垂直于直线 \( x + 2y + 1 = 0 \). (1) 求 a, b 的值; (2) 若函数 \( g(x) = \frac{\e^x}{f(x)} \)，讨论 \( g(x) \) 的单调性.
\end{problem}

\begin{problem}
如图，在四棱锥 \(S - ABCD\) 中，\(AD // BC\) 且 \(AD \perp CD\)；平面 \(CSD \perp\) 平面 \(ABCD, CS \perp DS, CS = 2AD = 2, E\) 为 \(BS\) 的中点，\(CE = \sqrt{2}\)，\(AS = \sqrt{3}\). 求： (1) 点 A 到平面 BCS 的距离; (2) 二面角 \(E - CD = A\) 的大小.
\end{problem}

\begin{problem}
已知以原点 \(O\) 为中心的椭圆的一条准线方程为 \(x = \frac{4\sqrt{3}}{3}\)，离心率 \(e = \frac{\sqrt{3}}{2}\)，则 \(S_{\triangle ABC}\) 的端点 \(A\) 在 \(AB\) 上的动点. (1) 若点 \(C, D\) 的坐标分别是 \((0, -\sqrt{3}), (0, \sqrt{3})\)，求 \(|MC| \cdot |MD|\) 的最大值； (2) 如图，点 \(A\) 的坐标为 \((1,0)\)，\(B\) 是圆 \(x^{2} + y^{2} = 1\) 上的点，\(N\) 是点 \(M\) 在 \(x\) 轴上的射影，点 \(Q\) 满足条件: \(\overrightarrow{OQ} = \overrightarrow{OM} + \overrightarrow{ON}\)，\(\overrightarrow{QA} \cdot \overrightarrow{BA} = 0\). 求线段 \(QB\) 的中点 \(P\) 的轨迹方程.
\end{problem}

\begin{problem}
设 \(m\) 个不全相等的正数 \(a_1, a_2, \cdots, a_m (m \geqslant 7)\) 依次围成一个圆圈. (1) 若 \(m = 2009\)，且 \(a_1, a_2, \cdots, a_{1005}\) 是公差为 \(d\) 的等差数列，则 \(a_1, a_{2009}, a_{2008}, \cdots, a_{1008}\) 是公比为 \(q\) 的等比数列; 数列 \(a_1, a_2, \cdots, a_m\) 的前 \(n\) 项和 \(S_n\) (\(n \leqslant m\)) 满足: \(S_1 = 15\)，\(S_{2009} = S_{2007} + 12a_1\)，求通项 \(a_n\) (\(n \leqslant m\)); (2) 若每个数 \(a_{n}\) (\(n \leqslant m\)) 是其左右相邻两数字的平方的等比中项，求证: \(a_{1} + \cdots + a_{n} + a_{1}^{2} + \cdots + a_{m}^{2} > mn; a_{1}^{2} \cdot a_{m}\).
\end{problem}

